\documentclass[student]{ne630boardhandout}
\hypersetup{
  pdftitle={NE 630 -- Lesson 07: Scattering Kinematics Derivations},
  pdfauthor={Jeremy Roberts},
  pdfsubject={Derivation supplement}
}

\providecommand{\muc}{\mu_{\mathrm c}}
\providecommand{\mul}{\mu_{\mathrm L}}
\providecommand{\thetac}{\theta_{\mathrm c}}
\providecommand{\thetaL}{\theta_{\mathrm L}}
\providecommand{\vc}{v_{\mathrm c}}
\providecommand{\bvc}{\mathbf{v}_{\mathrm c}}
\providecommand{\bVCM}{\mathbf{V}_{\mathrm{CM}}}
\providecommand{\VCM}{V_{\mathrm{CM}}}
\providecommand{\SourceEq}[1]{\par\nobreak\hfill{\sffamily\scriptsize\color{KSUGray}#1}\par}

\setlength{\abovedisplayskip}{4pt plus 1pt minus 1pt}
\setlength{\belowdisplayskip}{4pt plus 1pt minus 1pt}
\setlength{\abovedisplayshortskip}{3pt plus 1pt}
\setlength{\belowdisplayshortskip}{3pt plus 1pt}
\setlength{\jot}{2pt}

\HandoutSetup
  {NE 630}
  {07}
  {Scattering Kinematics: Derivations}
  {Textbook \S2.5; Rydin \S\S2.1--2.2.}

\begin{document}
{\sffamily\bfseries\large\color{KSUDeepPurple}
  Appendix: derivation of the kinematics equations\par}
\vspace{4pt}

\begin{handoutcolumns}
\HandoutSection{A0}{Set-up}
The target is initially at rest in the lab. Take $\hat{\mathbf x}$ along
the incident neutron velocity, so $\mathbf v=v\hat{\mathbf x}$.
Let $m_n$ be the neutron mass and $M_N$ the target mass. We use
$A\equiv M_N/m_n\simeq$ mass number with $A\geq1$.
\begin{align*}
 \bVCM&=\frac{m_n}{m_n+M_N}\mathbf v
       =\frac{1}{1+A}\mathbf v,\\
 \vc&=\|\mathbf v-\bVCM\|=\frac{A}{1+A}v.
\end{align*}
In the CM frame, elastic scattering preserves speeds and merely rotates
directions. In the scattering plane,
\[
\bvc'=\vc\big(\cos\thetac\,\hat{\mathbf x}
                  +\sin\thetac\,\hat{\mathbf y}\big),
\]
and the lab velocity is $\mathbf v'=\bvc'+\bVCM$.
Throughout, $\muc=\cos\thetac$ and $\mul=\cos\thetaL$.

\HandoutSection{A1}{Energy--angle link}
Using the square of the velocity magnitude,
\begin{align*}
 \frac{E'}{E}
  &=\frac{v'^2}{v^2}
   =\frac{\|\bvc'+\bVCM\|^2}{v^2}\\
  &=\frac{\vc^2+\VCM^2+2\vc\VCM\muc}{v^2}.
\end{align*}
Insert $\vc=Av/(1+A)$ and $\VCM=v/(1+A)$ to obtain
\begin{fixedbox}{Energy as a function of CM angle}
\[
\frac{E'}{E}=\frac{A^2+2A\muc+1}{(A+1)^2}.
\]
\SourceEq{Rydin 2.18}
\end{fixedbox}

\HandoutSection{A2}{Lab--CM angle link}
The lab scattering cosine is
\[
\mul=\frac{\mathbf v'\!\cdot\!\hat{\mathbf x}}{\|\mathbf v'\|}
  =\frac{\vc\muc+\VCM}{v'}.
\]
Divide numerator and denominator by $v$ and use the energy--angle result:
\[
\frac{v'}{v}=\sqrt{\frac{E'}{E}}
 =\frac{\sqrt{A^2+2A\muc+1}}{A+1}.
\]
Therefore,
\begin{fixedbox}{Lab angle as a function of CM angle}
\[
\mul=\frac{1+A\muc}{\sqrt{A^2+2A\muc+1}}.
\]
\SourceEq{Rydin 2.21}
\end{fixedbox}

\begin{boardbox}{Notes: velocity transformations}
\RuledLines{4}
\end{boardbox}

\switchcolumn
\RaggedRight
\HandoutSection{A3}{Average lab cosine}
For \textbf{isotropic CM scattering}, $\muc=\cos\thetac$ is uniform on $[-1,1]$.
Consequently,
\[
\overline\mu_{\mathrm L}
 =\frac12\int_{-1}^{1}\frac{1+A\muc}{\sqrt{A^2+2A\muc+1}}\,d\muc.
\]
Introduce the integration variable
\[
z=\sqrt{A^2+2A\muc+1}.
\]
Then
\[
d\muc=\frac{z}{A}\,dz,\qquad
 1+A\muc=\frac12\big[z^2-(A^2-1)\big].
\]
For the $A\geq1$ model, the integration limits are
\[
z\big|_{\muc=-1}=A-1,\qquad
  z\big|_{\muc=1}=A+1.
\]
Thus,
\begin{align*}
 \overline\mu_{\mathrm L}
  &=\frac12\,\frac{1}{2A}
    \int_{A-1}^{A+1}\big[z^2-(A^2-1)\big],dz\\[3pt]
  &=\frac{1}{4A}\left[\frac{z^3}{3}-(A^2-1)z\right]_{A-1}^{A+1}.
\end{align*}
Evaluating the limits gives
\begin{fixedbox}{Isotropic-in-CM average}
\[
\overline\mu_{\mathrm L}=\frac{2}{3A}.
\]
\SourceEq{Rydin 2.56}
\end{fixedbox}

\begin{boardbox}{Notes: substitution and limits}
\RuledLines{11}
\end{boardbox}

\begin{takeawaybox}[Reading and equation locators]
The Rydin and FNRP equation locators are retained from the lesson slides.
Assigned reading: textbook \S2.5 and Rydin,
\textit{Nuclear Reactor Theory and Design}, \S\S2.1--2.2.
\end{takeawaybox}
\end{handoutcolumns}

\end{document}
